\input{_preamble}
\begin{document}
\ex.\label{r:one}
\a.\sublabel{r:a} First.
\b. Second.
\c.\sublabel{r:c} Third.

RANGE \refrange{r:a}{r:c} PLAINREF \ref{r:one} PREF \pref{r:c}
LAST \Last{} NEXT \Next{}

%% \sublabel must record the label of the level it is used at.  Recording the
%% letter counter unconditionally makes every roman sub-sub-example under one
%% letter store the same value, and the range below then ends in that letter
%% ("1b-i--b") instead of in the roman numeral -- a silently wrong reference.
\ex.\label{r:two}
\a. Letter a of the second example.
\b.\sublabel{r:bletter} Letter b, which encloses the romans.
\a.\sublabel{r:ri} Roman one.
\b. Roman two.
\c.\sublabel{r:riii} Roman three.
\z.\z.

ROMANRANGE \refrange{r:ri}{r:riii} ROMANREF \ref{r:riii}
LETTERRANGE \refrange{r:bletter}{r:bletter}

\ex. Target of Next.

Footnote test\footnote{FN
\ex. First footnote example.

\ex. Second one.

FNLAST \Last{} FNLLAST \LLast{}} done.
\end{document}
